% !TEX root = ../hw02.tex
\leavevmode
\begin{enumerate}
  \item \textit{Boundedness.}
    For $k\geq2$, each of the $k-1$ factors $2,3,\ldots,k$
    is at least $2$, so
    \[
      k!=1\cdot2\cdots k\geq2^{k-1}.
    \]
    Consequently, for $n\geq2$,
    \[
      \begin{aligned}
        2\leq x_n
          &=2+\sum_{k=2}^n\frac1{k!}\\
          &\leq2+\sum_{k=2}^n\frac1{2^{k-1}}\\
          &=2+\frac{\tfrac12(1-(\tfrac12)^{n-1})}{1-\tfrac12}
          =3-2^{1-n}<3.
      \end{aligned}
    \]
    Since $x_1=2$, it follows that $2\leq x_n<3$ for every
    $n\geq1$. Thus $S$ is bounded.

  \item \textit{Infimum and supremum.}
    We take the infimum and supremum in $\R$. Since $x_1=2$
    and all terms are at least $2$, we have $a=\inf S=2$.
    Moreover, $x_{n+1}-x_n=1/(n+1)!>0$, so $(x_n)$ is strictly
    increasing. Since part (a) gives an
    upper bound, the monotone convergence theorem implies that
    $x_n\to L$ for some $L\in\R$.\sidenote{\textit{Monotone convergence theorem.}
      Every increasing real sequence that is bounded above converges.}
    Each $x_n\leq L$. If $u<L$, convergence gives a term $x_n>u$,
    so $u$ is not an upper bound. Therefore $b=\sup S=L$.

    All derivatives of $\exp(x)$ are continuous and equal to
    $\exp(x)$, so their values at $0$ are $1$.
    Taylor's theorem at $0$, evaluated at $1$, therefore
    gives\sidenote{\textit{Taylor remainder.} For a function with
      $n+1$ continuous derivatives on $[0,1]$, its degree-$n$
      expansion about $0$, evaluated at $1$, has remainder
      $f^{(n+1)}(\xi)/(n+1)!$ for some $\xi\in(0,1)$.}
    \[
      e=\sum_{k=0}^n\frac1{k!}
        +\frac{\exp(\xi_n)}{(n+1)!},\qquad 0<\xi_n<1.
    \]
    The remainder satisfies
    \[
      0<e-x_n\leq\frac{e}{(n+1)!}\longrightarrow0.
    \]
    Hence $L=e$, and
    \[
      a=2,\qquad b=e=\sum_{k=0}^{\infty}\frac1{k!}.
    \]

  \item \textit{Rationality.}
    We have $a=2\in\Q$, but $b=e\notin\Q$. To prove the latter,
    suppose that $e=p/q$ with $p,q\in\N$. Choose an integer
    $N\geq\max\{q,2\}$. Then
    \[
      R=N!\left(e-\sum_{k=0}^N\frac1{k!}\right)
    \]
    is an integer because both $N!e=N!p/q$ and every $N!/k!$ with
    $0\leq k\leq N$ are integers. On the other hand,
    \[
      \begin{aligned}
        0<R
          &=\sum_{j=1}^{\infty}
            \frac1{(N+1)(N+2)\cdots(N+j)}\\
          &\leq\sum_{j=1}^{\infty}\frac1{(N+1)^j}
          =\frac1N\leq\frac12<1.
      \end{aligned}
    \]
    This contradicts $R\in\Z$. Thus $e$ is irrational.

    In particular, $S$ has no supremum in the ordered field $\Q$.
    Indeed, any rational upper bound $u$ must satisfy $u>e$.
    By the density of $\Q$ in $\R$, there is a rational $v$
    with $e<v<u$, and $v$ is a smaller upper bound for $S$.
\end{enumerate}
